\chapter{理解大型語言模型}
\label{ch:understanding-llms}

% 章節摘要
\begin{quote}
本章介紹大型語言模型（LLM）的基本概念，包括它們的運作原理、應用場景，以及為什麼要從零開始打造自己的 LLM。
\end{quote}

\section{什麼是大型語言模型？}

\term{大型語言模型}{Large Language Model, LLM} 是一種基於深度學習的人工智慧模型，專門用於理解和生成自然語言文字。這些模型透過在大量文字資料上進行訓練，學習語言的模式、結構和語義。

\subsection{LLM 的關鍵特徵}

大型語言模型具有以下關鍵特徵：

\begin{itemize}
    \item \textbf{規模龐大}：通常包含數億到數千億個參數
    \item \textbf{自監督學習}：從未標記的文字資料中學習
    \item \textbf{情境學習}：能夠根據上下文理解和生成文字
    \item \textbf{少樣本學習}：僅需少量範例即可執行新任務
    \item \textbf{泛用性}：可應用於多種不同的自然語言處理任務
\end{itemize}

\subsection{LLM 的演進歷程}

圖 \ref{fig:llm-timeline} 展示了大型語言模型的發展歷程。

\begin{figure}[H]
    \centering
    % 此處將包含時間軸圖片
    % \includegraphics[width=0.9\textwidth]{ch01/llm-timeline.pdf}
    \caption{大型語言模型的發展時間軸}
    \label{fig:llm-timeline}
\end{figure}

\section{LLM 的運作原理}

\subsection{基本架構}

現代 LLM 大多基於 \term{Transformer}{Transformer} 架構，這是一種於 2017 年提出的深度學習架構。Transformer 的核心是 \term{注意力機制}{Attention Mechanism}，它能讓模型關注輸入序列中的重要部分。

\begin{definition}[Transformer 架構]
Transformer 是一種序列到序列模型架構，主要由以下元件組成：
\begin{itemize}
    \item 自注意力層（Self-Attention Layer）
    \item 前饋神經網路（Feed-Forward Network）
    \item 層正規化（Layer Normalization）
    \item 殘差連接（Residual Connection）
\end{itemize}
\end{definition}

\subsection{訓練過程}

LLM 的訓練通常分為兩個階段：

\begin{enumerate}
    \item \textbf{預訓練（Pre-training）}：
    \begin{itemize}
        \item 在大量未標記文字上訓練
        \item 學習語言的一般模式和知識
        \item 使用自監督學習目標，如下一標記預測
    \end{itemize}

    \item \textbf{微調（Fine-tuning）}：
    \begin{itemize}
        \item 在特定任務的資料上進一步訓練
        \item 調整模型以適應特定應用
        \item 可能包括指令微調或基於人類回饋的強化學習
    \end{itemize}
\end{enumerate}

\section{為什麼要從零開始打造 LLM？}

\subsection{學習目標}

從零開始實作 LLM 有以下好處：

\begin{itemize}
    \item \textbf{深入理解}：透過實作理解每個元件的運作方式
    \item \textbf{靈活性}：能夠根據需求自訂模型架構
    \item \textbf{除錯能力}：更容易診斷和解決問題
    \item \textbf{最佳化}：可針對特定用例進行效能最佳化
\end{itemize}

\begin{example}[教育用 LLM]
本書將引導您建立一個小型但功能完整的 GPT 模型，包含：
\begin{itemize}
    \item 約 124M 參數（與 GPT-2 small 相當）
    \item 可在一般筆記型電腦上訓練
    \item 能夠生成連貫的文字
    \item 可載入預訓練權重進行微調
\end{itemize}
\end{example}

\subsection{實用考量}

雖然訓練大型的生產級 LLM 需要龐大的計算資源，但從零開始打造較小的模型仍然非常實用：

\begin{itemize}
    \item 可用於特定領域的應用
    \item 適合資源受限的環境
    \item 便於實驗和研究新的架構
    \item 有助於理解和評估現有的大型模型
\end{itemize}

\section{本書概覽}

本書將帶領您完成以下學習旅程：

\subsection{章節架構}

\begin{description}
    \item[第 2 章] 處理文字資料 --- 學習如何將文字轉換為模型可處理的格式
    \item[第 3 章] 編寫注意力機制 --- 實作 Transformer 的核心元件
    \item[第 4 章] 從零開始實作 GPT 模型 --- 組合各個元件建立完整模型
    \item[第 5 章] 使用未標記資料進行預訓練 --- 訓練模型學習語言模式
    \item[第 6 章] 針對文字分類進行微調 --- 將模型應用於特定任務
    \item[第 7 章] 針對遵循指令進行微調 --- 訓練模型理解和執行指令
\end{description}

\subsection{學習路徑}

圖 \ref{fig:mental-model} 展示了本書涵蓋的內容架構。

\begin{figure}[H]
    \centering
    % \includegraphics[width=0.8\textwidth]{ch01/mental-model.jpg}
    \caption{LLM 開發的心智模型}
    \label{fig:mental-model}
\end{figure}

\section{先備知識}

\subsection{必要的背景知識}

要充分利用本書，您需要具備：

\begin{itemize}
    \item \textbf{Python 程式設計}：紮實的 Python 基礎
    \item \textbf{基礎數學}：線性代數、微積分和機率論的基本概念
    \item \textbf{（可選）深度學習}：有神經網路的基礎會有幫助
\end{itemize}

\subsection{PyTorch 基礎}

本書使用 PyTorch 作為深度學習框架。如果您不熟悉 PyTorch，附錄 A 提供了快速入門指南。

\begin{remark}
即使您是 PyTorch 新手，也可以跟隨本書學習。所有程式碼都有詳細的解釋，並且逐步建構。
\end{remark}

\section{開發環境設定}

\subsection{硬體需求}

本書的程式碼設計為可在一般筆記型電腦上執行：

\begin{itemize}
    \item \textbf{CPU}：現代多核心處理器
    \item \textbf{RAM}：至少 8GB（建議 16GB）
    \item \textbf{GPU}：可選，但會大幅加速訓練
    \item \textbf{儲存空間}：至少 10GB 可用空間
\end{itemize}

\subsection{軟體需求}

您需要安裝以下軟體：

\begin{itemize}
    \item Python 3.10 或更高版本
    \item PyTorch
    \item 其他 Python 套件（詳見 \texttt{requirements.txt}）
\end{itemize}

詳細的安裝指南請參閱 \texttt{setup} 目錄。

\section{本章小結}

本章介紹了：

\begin{itemize}
    \item 大型語言模型的基本概念和特徵
    \item LLM 的運作原理和訓練過程
    \item 從零開始實作 LLM 的價值和意義
    \item 本書的架構和學習路徑
    \item 必要的先備知識和開發環境設定
\end{itemize}

在下一章中，我們將開始處理文字資料，學習如何將原始文字轉換為模型可以理解的數值表示。

\subsection{延伸閱讀}

\begin{itemize}
    \item Vaswani et al., "Attention is All You Need" (2017)
    \item Radford et al., "Language Models are Unsupervised Multitask Learners" (2019)
    \item Brown et al., "Language Models are Few-Shot Learners" (2020)
\end{itemize}

\subsection{練習題}

本章無程式碼練習。建議：

\begin{enumerate}
    \item 研究並比較不同的大型語言模型（如 GPT、BERT、T5）
    \item 探索 LLM 的應用案例
    \item 設定您的開發環境並確認可以執行 PyTorch 程式碼
\end{enumerate}
